\documentclass[12pt,a4paper]{article}
\usepackage[top=16mm,bottom=18mm,left=16mm,right=16mm,headheight=0pt,headsep=0mm,footskip=10mm]{geometry}
\usepackage{iftex}
\ifPDFTeX
  \usepackage[utf8]{inputenc}
  \usepackage[T5]{fontenc}
  \usepackage{mathptmx}
\else
  \usepackage{fontspec}
  \usepackage{mathptmx}
  \IfFontExistsTF{Times New Roman}{\setmainfont{Times New Roman}}{\IfFontExistsTF{Tinos}{\setmainfont{Tinos}}{\setmainfont{TeX Gyre Termes}}}
\fi
\usepackage{amsmath,amssymb}
\usepackage{xcolor,tikz,graphicx,array,tabularx,enumitem,fancyhdr,lastpage}

% Bo cuc, font, mau so cau va khung khoanh dap an dua tren BTVN_Luong_giac_06.tex.
% Noi dung: De TDM11X33 -- Duong thang va mat phang song song (13 cau trac nghiem).
% Khong tao khoang trong de viet loi giai giua cac cau.
\definecolor{questionpurple}{RGB}{137,38,143}
\definecolor{answerblue}{RGB}{35,45,140}
\setlength{\parindent}{0pt}
\setlength{\parskip}{0pt}
\renewcommand{\baselinestretch}{1.08}
\pagestyle{fancy}
\fancyhf{}
\renewcommand{\headrulewidth}{0pt}
\fancyfoot[L]{\small Toán 11 -- Hình học không gian}
\fancyfoot[R]{\small Trang \thepage/\pageref{LastPage} -- Đề test số 08.1}
\renewcommand{\footrulewidth}{.4pt}

\newlist{questions}{enumerate}{1}
\setlist[questions]{label=\textcolor{questionpurple}{\bfseries Câu \arabic*.},
 leftmargin=0pt,itemindent=16mm,labelwidth=14mm,labelsep=2mm,align=left,
 itemsep=0pt,topsep=0pt,parsep=0pt,partopsep=0pt}
\newcommand{\answerlabel}[1]{\tikz[baseline=(choice.base)]\node[
 draw=black!65,text=answerblue,circle,inner sep=0pt,minimum size=4.8mm,
 line width=.4pt,font=\normalsize] (choice) {#1};}
\newcommand{\fourchoices}[4]{\par\begingroup\small
 \setlength{\tabcolsep}{0pt}\renewcommand{\arraystretch}{1.22}
 \begin{tabularx}{\linewidth}{@{}*{4}{>{\raggedright\arraybackslash}X}@{}}
 \answerlabel{A}~#1 & \answerlabel{B}~#2 &
 \answerlabel{C}~#3 & \answerlabel{D}~#4
 \end{tabularx}\endgroup}
\newcommand{\twochoices}[4]{\par\begingroup\small
 \setlength{\tabcolsep}{2pt}\renewcommand{\arraystretch}{1.25}
 \begin{tabularx}{\linewidth}{@{}*{2}{>{\raggedright\arraybackslash}X}@{}}
 \answerlabel{A}~#1 & \answerlabel{B}~#2\\
 \answerlabel{C}~#3 & \answerlabel{D}~#4
 \end{tabularx}\endgroup}
\newcommand{\onechoices}[4]{\par\begingroup\small
 \setlength{\tabcolsep}{1pt}\renewcommand{\arraystretch}{1.15}
 \begin{tabularx}{\linewidth}{@{}>{\raggedright\arraybackslash}X@{}}
 \answerlabel{A}~#1\\
 \answerlabel{B}~#2\\
 \answerlabel{C}~#3\\
 \answerlabel{D}~#4
 \end{tabularx}\endgroup}
\newcommand{\tdm}[1]{{\small\textcolor{answerblue}{[#1]}}}

\begin{document}
\noindent
\begin{minipage}[t]{.35\textwidth}
  \centering
  {\bfseries CHUYÊN ĐỀ -- TOÁN -- 11\par}
  \vspace{2pt}\rule{.72\linewidth}{.4pt}\par\vspace{3pt}
  HÌNH HỌC KHÔNG GIAN\par\vspace{4pt}
  {\itshape ĐỀ TEST SỐ 08.1\par}
\end{minipage}\hfill
\begin{minipage}[t]{.62\textwidth}
  \centering
  {\bfseries ĐƯỜNG THẲNG VÀ MẶT PHẲNG\par
  SONG SONG\par}
  \vspace{5pt}
  {\small\bfseries ĐỀ TỰ LUYỆN TDM11X33\par
  13 CÂU TRẮC NGHIỆM -- 75 PHÚT\par}
\end{minipage}
\vspace{4mm}
\noindent
\begin{minipage}[t]{.72\textwidth}
  \textbf{Họ, tên học sinh:}\dotfill\par
  \vspace{3pt}
  \textbf{Lớp:}\dotfill
\end{minipage}\hfill
\begin{minipage}[t]{.25\textwidth}
  \raggedleft\setlength{\fboxsep}{4pt}
  \fbox{\textbf{Đề test số: 08.1}}
\end{minipage}

\vspace{3mm}
\textbf{PHẦN A. CÁC BÀI TẬP NỀN TẢNG VÀ VẬN DỤNG TRỌNG ĐIỂM (13 CÂU).}
\begin{questions}
\item \tdm{TDM21} Mệnh đề \textbf{đúng} là
\onechoices
  {Đường thẳng và mặt phẳng song song với nhau nếu chúng không cắt nhau.}
  {Nếu hai đường thẳng cùng song song với một mặt phẳng thì chúng song song với nhau.}
  {Nếu hai đường thẳng cùng song song với một mặt phẳng thì chúng cắt nhau.}
  {Nếu đường thẳng $d$ không song song và không cắt mặt phẳng $(\alpha)$ thì $d$ nằm trong $(\alpha)$.}

\item \tdm{TDM21} Cho hai mặt phẳng $(\alpha)$, $(\beta)$ và hai đường thẳng $a,b$. Đáp án \textbf{sai} là
\onechoices
  {Nếu $a$ và $(\alpha)$ có hai điểm chung trở lên thì $(\alpha)$ chứa $a$.}
  {Nếu $b\subset(\alpha)$, $a\parallel b$ thì $a\parallel(\alpha)$.}
  {Nếu $a\parallel(\alpha)$, $a\subset(\beta)$, $(\alpha)\cap(\beta)=b$ thì $a\parallel b$.}
  {Nếu $(\alpha)\cap(\beta)=b$, $a\parallel(\alpha)$, $a\parallel(\beta)$ thì $a\parallel b$.}

\item \tdm{TDM22} Cho hai đường thẳng $a,b$. Đáp án \textbf{sai} là
\onechoices
  {Nếu $a$ và $b$ cắt nhau thì qua $a$ không tồn tại mặt phẳng nào song song với $b$.}
  {Nếu $a$ và $b$ trùng nhau thì qua $a$ không tồn tại mặt phẳng nào song song với $b$.}
  {Nếu $a$ và $b$ chéo nhau thì qua $a$ tồn tại vô số mặt phẳng song song với $b$.}
  {Nếu $a$ và $b$ song song với nhau thì qua $a$ tồn tại vô số mặt phẳng song song với $b$.}

\item \tdm{TDM21} Cho hình chóp $S.ABCD$ có đáy $ABCD$ là hình bình hành. Đáp án \textbf{sai} là:
\fourchoices
  {$AB\parallel(SCD)$.}
  {$AC\parallel(SBD)$.}
  {$AD\parallel(SBC)$.}
  {$CD\parallel(SAB)$.}

\item \tdm{TDM21} Cho tứ diện $ABCD$. Gọi $M$ là điểm nằm trên cạnh $AB$ của tứ diện $ABCD$. Mặt phẳng $(\alpha)$ đi qua $M$ và song song với $BD$. Gọi $N=AD\cap(\alpha)$. Đáp án \textbf{đúng} là
\fourchoices
  {$MN\parallel BD$.}
  {$MN\parallel CD$.}
  {$MN\parallel BC$.}
  {$MN\parallel AC$.}

\item \tdm{TDM32} Cho tứ diện $ABCD$. Gọi $M$ là trung điểm của cạnh $AB$ và $N$ là một điểm thuộc cạnh $CD$ không trùng với $C$ và $D$. Mặt phẳng $(\alpha)$ đi qua $MN$ và song song với $BC$. Gọi $P$ và $Q$ lần lượt là giao của $AC$ và $BD$ với $(\alpha)$. Biết rằng tứ giác $MPNQ$ là hình bình hành. Tỉ số $\dfrac{NC}{ND}$ bằng
\fourchoices
  {$\dfrac23$.}
  {$2$.}
  {$1$.}
  {$\dfrac12$.}

\item \tdm{TDM32} Cho hai hình bình hành $ABCD$ và $ABEF$ không cùng nằm trong một mặt phẳng. Gọi $P,Q$ lần lượt là trọng tâm của hai tam giác $ABF$ và $ABC$. Hỏi $PQ$ song song với mặt phẳng nào cho ở dưới đây?
\fourchoices
  {$(ACF)$.}
  {$(ACE)$.}
  {$(ABC)$.}
  {$(BCE)$.}

\item \tdm{TDM32} Cho hình chóp $S.ABCD$ có đáy $ABCD$ là hình bình hành. Gọi $M$ và $N$ lần lượt là trọng tâm của hai tam giác $ABC$ và $SAB$. Hỏi $MN$ song song với mặt phẳng nào cho ở dưới đây?
\fourchoices
  {$(SBD)$.}
  {$(SAC)$.}
  {$(SCD)$.}
  {$(SAD)$.}

\item \tdm{TDM42} Cho hình chóp $S.ABCD$ có đáy $ABCD$ là hình bình hành, có $SD=AB=6$. Gọi $M$ là trung điểm của cạnh $SA$. Mặt phẳng $(\alpha)$ đi qua $M$ song song với hai cạnh $AB$ và $SD$. Gọi $N,P,Q$ lần lượt là giao điểm của $AD,BC,SB$ với $(\alpha)$. Biết chu vi của tứ giác $MNPQ$ bằng $17$. Khi đó chu vi của tam giác $SCD$ bằng
\fourchoices
  {$20$.}
  {$24$.}
  {$18$.}
  {$22$.}

\item \tdm{TDM41} Cho tứ diện $ABCD$ có $AB+CD=10$. Gọi $M$ là trung điểm của $AB$. Mặt phẳng $(\alpha)$ đi qua $M$ song song với hai cạnh $AD$ và $BC$. Biết rằng $AC,CD,DB$ cắt $(\alpha)$ lần lượt tại $N,P,Q$. Chu vi của tứ giác $MNPQ$ bằng
\fourchoices
  {$15$.}
  {$12$.}
  {$10$.}
  {$20$.}

\item \tdm{TDM42} Cho hình chóp $S.ABCD$ có đáy $ABCD$ là hình bình hành. Gọi $M$ là trung điểm của cạnh $SC$. Mặt phẳng $(\alpha)$ qua $A$ và $M$ cắt các cạnh $SB$ và $SD$ lần lượt tại $B'$ và $D'$. Giá trị của biểu thức
$T=\dfrac{SB}{SB'}+\dfrac{SD}{SD'}$ bằng
\fourchoices
  {$2$.}
  {$4$.}
  {$3$.}
  {$\dfrac83$.}

\item \tdm{TDM52} Cho tứ diện $ABCD$. Gọi $G$ là trọng tâm của tứ diện và $M$ là một điểm nằm trên cạnh $BC$ sao cho $MB=2MC$. Mặt phẳng $(\alpha)$ đi qua $G$ và $M$ đồng thời song song với cạnh $AD$. Biết giao của đường thẳng $BD$ và $(\alpha)$ là $I$. Tỉ số $\dfrac{IB}{ID}$ bằng
\fourchoices
  {$3$.}
  {$4$.}
  {$\dfrac43$.}
  {$\dfrac72$.}

\item \tdm{TDM42} Cho hình chóp tứ giác $S.ABCD$ có đáy $ABCD$ là hình bình hành. Gọi $M$ là trung điểm của cạnh $SC$. Mặt phẳng $(\alpha)$ đi qua $AM$ và song song với $BD$. Gọi $P$ là giao điểm của cạnh $SB$ với mặt phẳng $(\alpha)$. Tỉ số $\dfrac{SP}{SB}$ bằng
\fourchoices
  {$\dfrac23$.}
  {$\dfrac13$.}
  {$\dfrac34$.}
  {$\dfrac12$.}

\end{questions}
\vspace{2mm}
\begin{center}\textbf{---------- HẾT ----------}\end{center}
\end{document}
